\documentclass[11pt,a4paper]{article}

% --- PACCHETTI FONDAMENTALI ---
\usepackage[utf8]{inputenc}
\usepackage[T1]{fontenc}
\usepackage[italian]{babel}
\usepackage{geometry}
\geometry{a4paper, margin=2.5cm}
\usepackage{amsmath, amssymb, amsthm}
\usepackage{graphicx}
\usepackage{booktabs}
\usepackage{hyperref}
\usepackage{xcolor}
\usepackage{listings}
\usepackage{fancyhdr}
\usepackage{caption}
\usepackage{subcaption}

% --- COLORI E CONFIGURAZIONE CODE HIGHLIGHTING ---
\definecolor{codegreen}{rgb}{0,0.6,0}
\definecolor{codegray}{rgb}{0.5,0.5,0.5}
\definecolor{codepurple}{rgb}{0.58,0,0.82}
\definecolor{backcolour}{rgb}{0.95,0.95,0.92}
\newcommand{\vir}[1]{``#1''}

\lstdefinestyle{mystyle}{
    backgroundcolor=\color{backcolour},
    commentstyle=\color{codegreen},
    keywordstyle=\color{magenta},
    numberstyle=\tiny\color{codegray},
    stringstyle=\color{codepurple},
    basicstyle=\ttfamily\small,
    breakatwhitespace=false,
    breaklines=true,
    captionpos=b,
    keepspaces=true,
    numbers=left,
    numbersep=5pt,
    showspaces=false,
    showstringspaces=false,
    showtabs=false,
    tabsize=2
}
\lstset{style=mystyle}

% --- Intestigazione HEADER E FOOTER ---
\pagestyle{fancy}
\fancyhf{}
\fancyhead[L]{Assignment 2}
\fancyhead[R]{Mattia Rey e Luca Ughetto}
\fancyfoot[C]{\thepage}

% --- DATI DEL DOCUMENTO ---
\title{
    \normalsize Università degli Studi di Torino \\
    \vspace{0.4cm}
    \Large \textbf{Corso di Human Computer Interaction} \\
    \vspace{0.3cm}
    \huge \textbf{Assignment 2}\\
    \Large \textbf{Gruppo DJ}
}
\author{
    \textbf{Autori:} Rey Mattia e Luca Ughetto \\
    \textbf{Matricole:}  1143424 e 1145888  \\
}
\date{2026}

\begin{document}

    \maketitle
    \thispagestyle{empty}

    \vspace{0.5cm}
    \begin{abstract}
        \noindent

    \end{abstract}

    \newpage
    \tableofcontents
    \newpage

% ==========================================
    \section{Introduzione}
    Giacomo
    % ==========================================

% =======================================
    \section{Definizione delle Personas}
% =======================================



\subsection{Persona 1: Marco Rossi (Progettista)}

\begin{table}[h!]
\centering
\small
\begin{tabular}{|p{3.5cm}|p{11.5cm}|}
\hline
\textbf{Nome} & Marco Rossi \\ \hline
\textbf{Età} & 34 anni \\ \hline
\textbf{Ruolo} & Coordinatore di Progetto presso un'associazione no-profit \\ \hline
\textbf{Aggettivo} & Pragmatico \\ \hline
\textbf{Citazione} & \textit{``Voglio trasformare le idee sociali della mia comunità in progetti concreti, ma la piattaforma non deve farmi perdere tempo con errori o navigazioni confuse.''} \\ \hline
\textbf{Piattaforme d'accesso} & Desktop/Laptop (60\%), Smartphone (40\%) \\ \hline
\end{tabular}
\caption{Profilo sintetico della Persona 1}
\label{tab:persona1}
\end{table}

\subsubsection{Giornata Tipo}
\begin{itemize}
    \item \textbf{07:30 -- 08:30:} Controllo rapido di e-mail e notifiche dallo smartphone durante il tragitto sui mezzi pubblici.
    \item \textbf{09:00 -- 13:00:} Coordinamento operativo delle attività dell'associazione e gestione dei volontari.
    \item \textbf{14:00 -- 17:00:} Lavoro al PC: stesura di proposte di progetto, gestione budget e preparazione materiali per il crowdfunding.
    \item \textbf{18:00 -- 20:00:} Gestione dei canali di comunicazione e risposte ai sostenitori.
    \item \textbf{21:00 -- 22:30:} Consultazione serale da laptop per pianificazione e ricerca di nuove opportunità formative.
\end{itemize}

\subsubsection{Bisogni Principali}
\begin{enumerate}
    \item **Trasparenza ed efficacia:** Raggiungere l'obiettivo economico stabilito tramite un canale affidabile.
    \item **Flusso di lavoro semplice:** Un processo chiaro e guidato per la creazione e pubblicazione della campagna.
    \item **Chiarezza informativa:** Accesso immediato a informazioni su formazione, partner e funzionamento del servizio.
    \item **Usabilità multi-dispositivo:** Esperienza coerente tra PC e smartphone.
\end{enumerate}

\subsubsection{Obiettivi di Miglioramento Personale}
\begin{itemize}
    \item Ridurre i tempi morti dovuti a problemi di usabilità dell'interfaccia.
    \item Massimizzare la conversione delle visite in donazioni grazie a un layout chiaro.
    \item Sviluppare competenze pratiche sul crowdfunding tramite sezioni formative integrate.
\end{itemize}

\subsubsection{Ostacoli e Pain Points}
\begin{itemize}
    \item \textbf{Sovraccarico cognitivo:} Home page caotica e ricca di informazioni superflue primarie.
    \item \textbf{Assenza di feedback:} Mancanza di indicatori di caricamento o messaggi di errore chiaro nella ricerca progetti.
    \item \textbf{Navigazione mobile limitata:} Menù non sticky e layout non sempre ottimizzato per smartphone.
    \item \textbf{Gestione errori nei form:} Moduli di registrazione senza validazione dinamica e perdita di dati inseriti in caso di errore.
    \item \textbf{Pagine non gestite:} Pagine di formazione prive di contenuti o elenchi partner con elementi non interattivi.
\end{itemize}


\end{document}